% ============================================================================
%  Chapter 4 -- Experimental Setup
% ----------------------------------------------------------------------------
%  FORMAT CONVERSION of thesis_writing/chapter_4/chapter_4_full_v1.md
%  (verified draft, written during the 2026-09-18 six-chapter structural
%  migration -- see that file's own provenance note). Scientific content,
%  wording, and every configuration value are preserved; only changes
%  required for valid LaTeX have been made (heading levels, math delimiters,
%  citation commands, cross-references, list/table environments). This
%  chapter states experimental SETTINGS only -- it reports no accuracy, BWT,
%  or forgetting number; those belong to Chapter 5 (Results), still a
%  placeholder as of this conversion.
%
%  LABELS: sec:ch4-* for sections, tab:ch4-* for tables, mirroring Chapter
%    3's sec:ch3-*/tab:ch3-* convention. Chapter label \label{ch:experimental-setup}
%    is unchanged from the pre-existing placeholder, so every \ref to it
%    already present in Chapters 1-3 continues to resolve.
%
%  CROSS-REFERENCES INTO CHAPTER 5: only the chapter-level \ref{ch:results}
%    is used; stale plain-text Chapter 5 subsection numbers were removed in
%    the 2026-09-23 cleanup pass.
%
%  BIBLIOGRAPHY: one new entry added, krizhevsky2009learning (CIFAR-100),
%    appended to references.bib. radford2021learning (CLIP) already existed,
%    reused from Chapter 3.
% ============================================================================

\chapter{Experimental Setup}
\label{ch:experimental-setup}

\section{Experimental Objectives}
\label{sec:ch4-objectives}

This chapter specifies the concrete experimental settings under which the
two continual-LoRA integration strategies defined in
Chapter~\ref{ch:methodology} --- SimpleAvg (Section~\ref{sec:ch3-simpleavg})
and RankExt (Section~\ref{sec:ch3-rankext}) --- together with the knowledge
distillation and FactorOrth stabilisation mechanisms of
Section~\ref{sec:ch3-stabilization}, were trained and evaluated. Three
objectives structure the chapter: to fix, once, the datasets, protocols,
backbone, and training configuration shared by the experiments reported in
Chapter~\ref{ch:results}, so results are not repeatedly re-described; to
state exactly which method variants and configuration differences exist
across the experiments this thesis draws on, since more than one
configuration of each strategy has been run; and to record the
implementation and reproducibility details --- software environment, seeds,
compute infrastructure, checkpointing --- needed to situate those results as
reproducible, traceable experiments rather than one-off numbers. This
chapter states settings only; Chapter~\ref{ch:results} reports and
interprets what those settings produced.

\section{Datasets}
\label{sec:ch4-datasets}

Two datasets are used. \textbf{CIFAR-100} \parencite{krizhevsky2009learning}
is the primary benchmark: 100 fine-grained object classes, 50,000 training
images and 10,000 test images, loaded in its native (unshuffled) label
order. Unless stated otherwise, every dataset-specific setting in this
chapter refers to CIFAR-100.

\textbf{CUB-200-2011} (Caltech-UCSD Birds-200-2011) provides a second-dataset
supporting evaluation. It has 200 bird-species classes and is loaded through
the dataset identifier \texttt{Donghyun99/\allowbreak CUB-200-2011}. The official split
is retained: the 5,994 images of the official training split are the
training source, and the 5,794 images of the official test split are used
only for final evaluation. The validation set is carved from the official
training split alone, with 5 images per class selected deterministically
under the fixed seed; the test split is never used for validation. The
classes are presented in a fixed seed-42 permutation rather than in the
native class order; this permutation is stored once and shared by every
method.

\section{Class-Incremental Protocol}
\label{sec:ch4-protocols}

The primary experimental setting used throughout this thesis is CIFAR-100
under a $5\times20$ class-incremental protocol, comprising five steps of 20
classes each. Unless explicitly stated otherwise, the main comparisons and
analyses in Chapter~\ref{ch:results} refer to this setting. It instantiates
the class-incremental setting of
Section~\ref{sec:ch3-problem-formulation} --- a fixed, disjoint,
non-revisited sequence of class groups, introduced under the native
(unshuffled) label order --- with class groups $[0,19]$, $[20,39]$,
$[40,59]$, $[60,79]$, $[80,99]$, and RankExt's cumulative rank schedule
$[16,32,48,64,80]$ (Section~\ref{sec:ch4-backbone}). This is the setting
used for the main, headline results in Chapter~\ref{ch:results}.

Every main CIFAR-100 result in Chapter~\ref{ch:results} uses the primary
$5\times20$ protocol. The only alternative CIFAR-100 partition reported is
a secondary $4\times25$ diagnostic for the RankExt family
(Section~\ref{sec:ch5-ablation}); a systematic study of protocol depth is
left to future work (Section~\ref{sec:ch6-future-work}).

CUB-200-2011 is evaluated under a $5\times40$ protocol: the 200 permuted
class indices are split into five steps of 40 classes each ($[0,39]$,
$[40,79]$, $[80,119]$, $[120,159]$, $[160,199]$). RankExt uses the same
cumulative rank schedule $[16,32,48,64,80]$ as on CIFAR-100, so each step
still adds a rank-16 block and the final cumulative rank is again 80, while
each step now introduces 40 rather than 20 classes; the schedule is
therefore not matched per class across the two datasets. No accuracy
number appears in this section; results are reported in
Chapter~\ref{ch:results}.

\section{Backbone and Adaptation Architecture}
\label{sec:ch4-backbone}

The backbone is the vision encoder of CLIP ViT-B/16
(\texttt{openai/clip-vit-base-patch16}; \parencite{radford2021learning}),
used identically for both integration strategies and frozen throughout
training (Section~\ref{sec:ch3-lora-framework}); CLIP's text encoder plays
no role. The pooled representation feeds a single linear classifier
whose output dimension equals the number of classes of the dataset:
$\mathbb{R}^{768}\to\mathbb{R}^{100}$ for CIFAR-100 and
$\mathbb{R}^{768}\to\mathbb{R}^{200}$ for CUB-200-2011
(Section~\ref{sec:ch3-classifier-construction}).
LoRA is injected into the query and value projection matrices
(\texttt{q\_proj}, \texttt{v\_proj}) of every attention block --- the same
two-module target set for both strategies; no alternative target-module set
was evaluated for any result in this thesis.

SimpleAvg is configured, throughout every experiment presented as primary in
Chapter~\ref{ch:results}, at rank $r=80$, $\alpha=160$, giving scaling
$s=\alpha/r=2$ (Section~\ref{sec:ch3-lora-framework}), fixed for every step.
One independent specialist is trained per step from the pretrained backbone;
after the final step, the reconstructed dense updates of the five
specialists are averaged arithmetically and added to a fresh pretrained
backbone, and each class's classifier row is copied from the specialist
that trained it (Section~\ref{sec:ch3-simpleavg}).
RankExt's primary $5\times20$ cumulative rank schedule is
$[16,32,48,64,80]$ (Section~\ref{sec:ch4-protocols} above), with a constant
per-rank scaling of 2.0 at every growth step, matching SimpleAvg's own
scaling by construction (Section~\ref{sec:ch3-lora-framework}). A single
persistent adaptation is carried across steps: at each step a new rank-16
block is appended, while all previously learned blocks remain active in the
forward pass and are hard-frozen (Section~\ref{sec:ch3-rankext}).

The CUB-200-2011 experiment uses the same backbone, target modules, method
definitions, and loss configuration as CIFAR-100
(Sections~\ref{sec:ch4-training}--\ref{sec:ch4-rankext-auxiliary}), with
the differences summarised in Table~\ref{tab:ch4-cub-config}. The most
important one is the RankExt classifier-head learning-rate multiplier,
which is $\times2$ in the refined CUB-200-2011 configuration and $\times1$
in the canonical CIFAR-100 benchmark. The CUB-200-2011 experiment is
therefore a second-dataset supporting evaluation of the same method
families, not a strictly controlled dataset-only replication of the
CIFAR-100 benchmark.

\begin{table}[htbp]
\centering
\begin{tabularx}{\textwidth}{@{}l X X@{}}
\toprule
\textbf{Setting} & \textbf{CIFAR-100 (primary)} & \textbf{CUB-200-2011} \\
\midrule
Classes / protocol & 100 / $5\times20$ & 200 / $5\times40$ \\
Class order & Native label order & Fixed seed-42 permutation \\
Classifier & $768\to100$ & $768\to200$ \\
Validation & 25 images per class, from the training set & 5 images per class, from the official training split \\
Final evaluation & CIFAR-100 test set & Official CUB-200-2011 test split \\
RankExt head-LR multiplier & $\times1$ & $\times2$ \\
Preprocessing & Resize, random crop (padding 8), horizontal flip, colour jitter & Aspect-ratio-preserving CLIP resize, random crop (padding 8), horizontal flip, light colour jitter; evaluation with deterministic centre crop \\
\bottomrule
\end{tabularx}
\caption[CIFAR-100 and CUB-200-2011 configuration differences]{Settings that
differ between the CIFAR-100 primary benchmark and the CUB-200-2011
evaluation. All other settings --- backbone, target modules, SimpleAvg rank
80, RankExt schedule $[16,32,48,64,80]$, learning rates, SimpleAvg head-LR
multiplier $\times10$, a maximum of 9 epochs per step, knowledge
distillation, FactorOrth, and the projected feature-protection term --- are
shared.}
\label{tab:ch4-cub-config}
\end{table}

\section{Training Configuration}
\label{sec:ch4-training}

The following training configuration is shared by every method variant and
every experiment presented as primary in Chapter~\ref{ch:results}, under the
primary $5\times20$ protocol:

\begin{table}[htbp]
\centering
\begin{tabularx}{\textwidth}{@{}l X@{}}
\toprule
\textbf{Setting} & \textbf{Value} \\
\midrule
Epochs per step & Maximum 9; best-validation checkpoint selection throughout. Plain SimpleAvg alone uses validation-based early stopping (patience 3) and stopped after 5--8 epochs per step; every other variant ran all 9 \\
Optimizer & AdamW \\
Learning rate & $5\times10^{-5}$ (SimpleAvg) / $1\times10^{-4}$ (RankExt) --- family-specific \\
Classifier-head LR multiplier & $\times10$ (SimpleAvg) / $\times1$ (RankExt) --- family-specific \\
Weight decay & 0.05 \\
LR scheduler & Cosine, with a 10\% warmup ratio \\
Batch size & 16 (effective batch 16, gradient accumulation 1) \\
Gradient clipping & Default Hugging Face \texttt{Trainer} behaviour, \texttt{max\_grad\_norm=1.0} \\
Mixed precision & fp16 when a GPU is available; fp32 otherwise \\
Replay & \textbf{None.} Every method variant in this thesis is rehearsal-free (Section~\ref{sec:ch3-problem-formulation}); no replay buffer is used by any reported method variant \\
Data augmentation (train) & Resize $\to$ random crop (padding 8) $\to$ random horizontal flip $\to$ colour jitter $\to$ normalise (CLIP statistics) \\
Data augmentation (eval) & Resize $\to$ normalise only, no augmentation \\
\bottomrule
\end{tabularx}
\caption[Training configuration]{Training configuration shared by every
method variant and every experiment presented as primary in
Chapter~\ref{ch:results}.}
\label{tab:ch4-training-config}
\end{table}

\textbf{Validation and best-epoch selection.} A held-out, class-stratified
validation split (25 images per class, deterministically shuffled) is
constructed once, before any method's training loop begins, and shared by
every method variant. At the end of each local epoch, pure validation
cross-entropy --- not the regularised training loss any KD- or
orthogonality-bearing method optimises --- is computed; whenever it improves on
the running best, the current trainable parameters are snapshotted, and the
best snapshot is reloaded once that step's training completes, before that
step's evaluation runs. This selection rule is identical across every method
variant and every experiment referenced in Chapter~\ref{ch:results}.

\textbf{Calibration.} The confidence-weighted, regime-grouped row-norm
calibration mode of Section~\ref{sec:ch3-calibration} is applied,
identically, to every method variant's final classifier before evaluation,
with confidence exponent $\gamma=0.65$ and per-step scaling factors clipped
to $[0.85, 1.30]$. As already established in
Section~\ref{sec:ch3-calibration}, this rescaling is post-hoc, touches only
the classifier's weight rows (the bias vector $b$ is never modified), and
the same calibrated classifier is used for both open and restricted
evaluation.

\subsection{Variant-Specific Loss Composition}
\label{sec:ch4-loss-composition}

Section~\ref{sec:ch3-stabilization} states that every variant's exact loss
composition, auxiliary stabilisation terms, and warmup schedules are
specified here. Each stabilisation mechanism's general form and the
temperature/weight symbols below are defined in
Sections~\ref{sec:ch3-kd} and~\ref{sec:ch3-factororth}; this section states
only the experimental configuration actually used, not a re-derivation of
either mechanism.

\textbf{Knowledge distillation.} Temperature $\tau=2$, weight $\kappa=1.0$,
with the $\tau^2$ scaling of Equation~\eqref{eq:ch3-kd}. Distillation is
active from the second class-incremental step onward (a teacher requires at
least one already-trained step) and is computed on the current step's own
input batch only. Its scope and schedule are family-specific. For RankExt,
it compares the full output distribution over all classes of the
classifier, which covers every class seen so far, and
there is no distillation-specific warmup: $\tau$ and $\kappa$ are constant
whenever distillation is active. For SimpleAvg, the comparison is
restricted to the previously seen classes, and the distillation weight is
ramped linearly from zero over the first epoch of each step. Teacher
construction is as defined in Section~\ref{sec:ch3-kd}: for SimpleAvg, a
frozen running merged prefix snapshot over the steps already trained; for
RankExt, the frozen persistent model as it stood at the end of the previous
step.

\textbf{FactorOrth.} The FactorOrth configuration is family-specific
(Section~\ref{sec:ch3-factororth}). SimpleAvg's FactorOrth variants use
dense-update orthogonality (Equation~\eqref{eq:ch3-factororth-simpleavg})
with $\lambda_{orth}=20$; RankExt's FactorOrth variants use factor-space
orthogonality on normalised factors
(Equations~\eqref{eq:ch3-factororth-a}--\eqref{eq:ch3-factororth-b}) with
$\lambda_{orth}=50$. A one-epoch linear warmup on the FactorOrth coefficient
is applied to all four FactorOrth-bearing variants: SimpleAvg + FactorOrth,
SimpleAvg + FactorOrth + KD, RankExt + FactorOrth, and RankExt + FactorOrth
+ KD + Protect. This
orthogonality-loss warmup is distinct from the $10\%$ learning-rate warmup
of the main optimizer schedule (Table~\ref{tab:ch4-training-config}).

\subsection{RankExt Auxiliary Stabilisation}
\label{sec:ch4-rankext-auxiliary}

Section~\ref{sec:ch3-eight-variants} notes that activating knowledge
distillation for RankExt also couples it to further, RankExt-specific
configuration choices, not present for SimpleAvg. These choices are
implementation/configuration details of the RankExt variants used in the
experiments, not a novel contribution of this thesis:

\begin{itemize}
\item Non-KD RankExt variants use a pretrained-backbone feature anchor
(weight $1.0$), a cosine-distance term against the frozen pretrained
backbone's own features.
\item KD-bearing RankExt variants instead use projected feature protection
(weight $30.0$), active from the second step onward, which discourages
drift of the shared representation along directions the earlier steps'
classifier rows already occupy.
\item The feature anchor and projected feature protection are mutually
exclusive: exactly one, or neither, is active for a given RankExt variant,
never both.
\item New-block output warmup (a one-epoch linear ramp on the
newly-appended rank block's contribution to the forward pass) is active for
non-KD RankExt variants and disabled for KD-bearing RankExt variants.
\end{itemize}

\section{Evaluation Metrics}
\label{sec:ch4-metrics}

Chapter~\ref{ch:methodology} defines every metric in this list precisely
(Sections~\ref{sec:ch3-open-restricted} and~\ref{sec:ch3-bwt-forgetting});
this section states only which are computed and reported, not their
derivation.

Three summary metrics are used throughout Chapter~\ref{ch:results} to
characterise the final model in a single comparison table; they are
distinct from the formal continual-learning metrics, backward transfer and
forgetting, defined next.

\begin{itemize}
\item \textbf{First-step accuracy} --- open accuracy on the first step's own
class group, at the final model.
\item \textbf{All-seen accuracy} --- open accuracy over every class
introduced by the point of evaluation, at the final model; the headline
comparison metric used throughout Chapter~\ref{ch:results}.
\item \textbf{Later-steps accuracy} --- open accuracy of the final model on
a single combined evaluation set formed by pooling every class introduced
after the first step; this is one accuracy number over the union of those
classes, not an average of separately computed per-step accuracies.
\end{itemize}

\begin{itemize}
\item \textbf{Open accuracy} and \textbf{restricted accuracy} --- computed
from the same model's logits on the same examples; open accuracy competes
every seen class, restricted accuracy masks to one step's own class group.
Both open and restricted class-group accuracies are reported using the
final calibrated model.
\item \textbf{Backward transfer (BWT)} and \textbf{forgetting} --- computed
exactly as defined in Section~\ref{sec:ch3-bwt-forgetting}. BWT is reported
for both integration strategies (with SimpleAvg's diagonal reference
substituted by its pre-merge specialist accuracy, since no genuine
intermediate checkpoint exists for that strategy); forgetting is reported
for RankExt only.
\item \textbf{Forward transfer} is not reported because it was not
validated sufficiently for reliable interpretation.
\end{itemize}

\section{Implementation and Reproducibility}
\label{sec:ch4-reproducibility}

\textbf{Software.} Training is implemented in PyTorch using the Hugging Face
\texttt{transformers.Trainer} API, the CLIP vision encoder from the Hugging
Face \texttt{transformers} model hub, and PEFT's
\texttt{LoraConfig}/\texttt{get\_peft\_model} for SimpleAvg's per-step
adapters (RankExt uses a custom \texttt{Growing\allowbreak RankLoRA\allowbreak Linear} module rather
than PEFT, to support persistent, growing-rank blocks ---
Section~\ref{sec:ch3-rankext}). The exact PyTorch/CUDA version used for any
specific run was not independently recovered from a saved environment log
during the preparation of this thesis and is not stated here rather than
guessed.

\textbf{Compute environment.} The CIFAR-100 experiments were performed on
the University of Padova DEI computing cluster using one NVIDIA A40 GPU per job, with jobs
submitted through SLURM.

\textbf{Seeds.} A single global seed drives every source of randomness
(Python's \texttt{random}, NumPy, and PyTorch) at process start, with an
additional per-method reseed inside the outer training loop so that no
method's initialisation or dropout draws depend on how many other methods
have already trained earlier in the same process. The seed is 42 for the
CIFAR-100 primary benchmark and for the CUB-200-2011 evaluation, where it
also determines the class permutation and the validation selection. Both
are single-seed results; neither is replicated across seeds.

\textbf{Checkpointing, logging, and result artifacts.} Each run writes, per
method and per class-incremental step: the best-validation-CE snapshot of
trainable parameters (Section~\ref{sec:ch4-training}), a per-epoch
training/validation loss log, and a structured set of result tables
(per-method final metrics, per-step open/restricted accuracy, best-epoch
selection records, and --- where applicable --- classifier-row-restoration
and rank-structure diagnostics) as CSV files under a per-run results
directory, together with a machine-readable configuration snapshot
recording the hyperparameters and configuration used for the run. All reported metrics
in Chapter~\ref{ch:results} are derived from these saved experiment
artifacts.
